\section{Exploration vs. Exploitation}
Learning agents needs to trade off two things:
\begin{itemize}
    \item \textbf{Exploration:} increase the amount of knowledge
    \item \textbf{Exploitation:} maximize the performance based on current knowledge
\end{itemize}

\subsection{K-armed Bandit}
Consider the following learning problem. You are faced repeatedly with a choice among $k$ options, or actions. After each choice you receive a numerical reward chosen from a stationary probability distribution that depends on the action you selected. Your objective is to maximize the expected total reward over some period, for example, over 1000 actions selections, or \textit{time steps}.
\begin{itemize}
    \item The original problem was named by an analogy to a slot machines with $k$ levers, where each action is like a play of one of the slot machine's levers and the rewards are the payoffs for hitting the jackpot.
\end{itemize}

\subsubsection{Greedy algorithms}\label{lab:greedy_algorithms}

In our problem each of the $k$ actions has an expected/mean reward given that that action is selected; let us call this the \textbf{value} of that action. The value of an arbitrary action $a$, denoted $q_\star (a)$, is the expected reward, given that $a$ is selected:
\[ 
    q_\star (a) = \mathbb{E}[R_t \mid A_t=a]
\]

Note that here the $q$ function does not depend on the state $S_t=s$, and that's because we have only one state in the K-armed bandit problem, so it's unnecessary to write it down.

We assume that we do not know the action values with certainty, although we have the estimated value of action $a$ at time step $t$, denoted as $Q_t(a)$. We'd like $Q_t(a)$ to be as close as possible to $q_\star (a)$. If we maintain estimates of the action values, then at any time step there is at least one action whose estimated value is greatest: we call these the \emph{greedy} actions. When we select one of these actions, we say that we are \textbf{exploiting} the current knowledge of the value of the actions. If, instead, we select one of the non-greedy actions then we say we are \textbf{exploring}, because this enables to improve the estimate of the non-greedy action values.


Recall that the true value of an action is the mean reward when that action is selected. A simple estimate of this quantity is the following:
\[  
    Q_t(a) = \frac{\text{sum of rewards when $a$ taken prior to $t$}}{\text{number of times $a$ taken prior to $t$}} = \frac{\sum_{i=1}^{t-1} R_i \cdot \mathbbm{1}(A_i=a)}{\sum_{i=1}^{t-1} \mathbbm{1}({A_i=a})}
\]

Where $\mathbbm{1}(predicate)$ denotes the random variable that is 1 if $predicate$ is true and 0 if it is not. Here the simplest action selection rule is to select one of the actions with the highest estimated value:
\[
    A_t = \argmax_{a} Q_t(a)
\]

That is a \textbf{greedy} action selection, that always exploits current knowledge to maximize immediate reward. Note the greediness stops the algorithm from sampling apparently inferior actions to see if they might really be better! That's why $\varepsilon$-$greedy$ methods got introduced: with a small probability $\varepsilon$ we select the action randomly from among the all possible actions. In the remaining cases we use normal greedy action selection.

\subsubsection{Gradient bandits}
Now the question is: can we learn policies $\pi(a)$ directly, instead of learning values? For instance, let's define an \textbf{action preference} function $H_t(a) \in \mathbb{R}$, a sort of "liking" for a certain action $a$. To turn these abstract preferences into actual probabilities, we use the \textbf{Softmax} distribution:
\[
    \mathbb{P}(A_t=a) \doteq \frac{e^{H_t(a)}}{\sum_{b} e^{H_t(b)}} \doteq \pi_t(a)
\]
Where $\pi_t(a)$ represents the probability of taking action $a$ at time $t$.

There is a natural learning algorithm for softmax action preferences based on the idea of stochastic gradient ascent. On each step, after selecting action $A_t$ and receiving the reward $R_t$, the action preferences are updated by:
\begin{equation}
    H_{t+1}(a) =
    \begin{cases}
        H_t(a) + \alpha R_t \cdot (1 - \pi_t(a)) &\qquad \text{if } a=A_t \\
        H_t(a) - \alpha R_t \cdot \pi_t(a) &\qquad \text{if } a \neq A_t
    \end{cases}
\end{equation}
This update mechanism creates a "push-pull" effect that ensures the policy remains a valid probability distribution:
\begin{itemize}
    \item For the chosen action $a=A_t$ we "push up" the preference. The term $(1-\pi_t(a))$ acts as a scaling factor, so that if the action was already likely the update is small, else the update is much larger
    \item For all other actions $a \neq A_t$ we "pull down" the preferences by subtracting $\alpha R_t \pi_t(a)$, we ensure that as one action becomes more desirable, the others become relatively less likely, maintaining the balance of the softmax distribution.
\end{itemize}

Now we can go ahead and see the gradient bandit algorithm as a stochastic approximation of a gradient ascent method; in \textit{exact} gradient descent, each action preference $H_t(a)$ would be incremented in proportion to the increment's effect on performance:
\[
    H_{t+1}(a) \doteq H_t(a) + \alpha \frac{\partial \mathbb{E}[R_t]}{\partial H_t(a)}
\]

with $\mathbb{E}[R_t] = \sum_x \pi_t(x) q_\star(x)$. Obviously this is impossible to implement, since we do not know the optimal $q_\star$ value.

\begin{tcolorbox}[colframe=black, coltitle=white, title=Function Approximation Primers, fonttitle=\bfseries]
    That's all good, but now how do we learn these preferences? Since we aren't just averaging rewards anymore, we update $H_t(a)$ using \textbf{Stochastic Gradient Ascent}. We now define the performance of our policy $\pi_\theta$ as the expected reward we get when following it. We call this $J(\theta)$: 

\[
    J(\theta) \doteq \mathbb{E}[R_t \mid \pi_\theta] = \sum_a \pi_\theta(a) \cdot q_\star(a)
\]
Our goal is to find the parameters $\theta$ (in our case, the preferences $H_t(a)$) that make $J(\theta)$ as large as possible. Since we want to maximize $J(\theta)$, we move our parameters in the direction that increases the expected reward most steeply. This direction is the \textbf{gradient}:
\[
    \theta_{t+1} = \theta_t + \alpha \nabla_\theta \mathbb{E}[R_t \, | \, \pi_{\theta_t}]
\]
Where:
\begin{itemize}
    \item $\theta_t$: our current internal preferences.
    \item $\alpha$: the step size
\end{itemize}

Now we would like to have $\nabla_\theta \mathbb{E}[R_t \mid\pi_\theta]$, but there's a problem: the expectation is over actions sampled from $\pi_\theta$, and $\theta$ lives inside that distribution, and we can't just push the gradient through a sampling operation the way we would with a normal function. The \textbf{log-likelihood trick} (also called REINFORCE) is the standard fix to this issue: it rewrites the gradient of an expectation as an expectation of a gradient

\[
    \nabla_\theta \mathbb{E}[R_t \mid \pi_\theta] = \mathbb{E}[R_t \nabla_\theta \ln (\pi_\theta(a))]
\]
Now, we need to find $\nabla_\theta \ln \pi_\theta(a)$ specifically for the Softmax function. In our case, the parameters $\theta = H_t(a)$. We need to calculate how the log-probability of an action a changes when we nudge a specific preference $H_t(k)$. Taking the partial derivative of the softmax gives us two different scenarios:
\begin{equation}
    \frac{\partial \ln \pi_t(a)}{\partial H_t(k)} = 
    \begin{cases} 
        1 - \pi_t(a) & \text{if } k = a \\
        -\pi_t(k)    & \text{if } k \neq a 
    \end{cases}
\end{equation}

Now let's imagine what happens if our rewards $R_t$ are large and positive: then the basic update will push all action preferences up significantly, even for the "bad" actions. The model eventually figures it out, but it's a slow, noisy process. The solution is to introduce a so-called \textbf{baseline} $b$ to the gradient-ascent rule update:

\[
    \theta_{t+1} = \theta_t + \alpha (R_t - b) \nabla_\theta \ln \pi_\theta (A_t)
\]

Mathematically everything works out as expected, since adding the baseline does not change the expected update:

$$\mathbb{E}_{a\sim\pi_\theta}\big[b\,\nabla_\theta \ln \pi_\theta(a)\big] =
b\sum_a \nabla_\theta \pi_\theta(a) = b\,\nabla_\theta\sum_a \pi_\theta(a) = b\,\nabla_\theta(1) = 0$$

since $b$ doesn't depend on the action being summed over, and probabilities always sum to $1$
\end{tcolorbox}



\subsection{How well can we do?}
While the Policy Gradient method with a baseline helps us converge faster, the Lai \& Robbins Theorem (1985) reminds us that exploration has a mandatory cost. Lai \& Robbins proved that for any reasonable strategy, the total regret must grow at least \textbf{logarithmically} with the number of steps $t$. Before going into the conclusion, we need to introduce some other concept needed for understanding it. First, let's remind ourselves what $q_\star$ is:

\[
    q_\star = \max_{a \in A} q(a) = \max_{a} \mathbb{E}[R_t \mid A_t = a]
\]
which is the \textbf{optimal value} achievable by any single action in the current environment. Instead, we call $q_\star(a)$ the \textbf{true value} (e.g. the expected reward) of a specific action $a$. We then define the \textbf{regret} of taking an action as:

\[
    \Delta_a \doteq q_\star - q_\star(a) 
\]
In other words, the regret represents the gap between the best action possible and the action $a$ which was actually chosen. The last piece we need to introduce is the \textbf{total regret} which, you can imagine, is the sum of all the single regrets:

\[
    L_t = \sum_{n=1}^{t} \left( q_\star - q(A_n) \right) = \sum_{n=1}^t \Delta_{A_n}
\]

Before stating Lai \& Robbins' result we need one tool from information theory: a way to measure "how different" two probability distribution are and that's exactly what KL divergence gives us. Intuitively: if pulling arm $a$ and pulling the optimal arm produce reward distributions that look almost identical, no amount of data will let you tell them apart quickly — you're forced to keep sampling the suboptimal arm just to gather enough evidence. That "forced sampling cost" is exactly what shows up as $1/D_{KL}$ in the regret bound below: small KL divergence $\Rightarrow$ big minimum regret.

\begin{definition}[Kullback-Leibler Divergence]
    The \textbf{Kullback-Leibler divergence} $D_{KL}(p \parallel q)$ is a measure from information theory that quantifies how much one probability distribution $p$ deviates from a second reference probability distribution $q$

    \[
        D_{KL}(p \parallel q) = \sum_{x \in \mathcal{X}} p(x) \log \frac{p(x)}{q(x)}
    \]
\end{definition}

\begin{theorem}[Lai \& Robbins]
    The asymptotic total regret is \textbf{at least logarithmic} in the number of steps $t$
    \[
        \liminf_{t\to\infty} \frac{L_t}{\ln t} \;\ge\; \sum_{a:\Delta_a>0} \frac{\Delta_a}{D_{KL}(p_a \,\|\, p_\star)}
    \]
\end{theorem}

\noindent
Let's dive into the formula a bit more:
\begin{itemize}
    \item We are summing over all $\Delta_a > 0$, that means we are summing over all suboptimal actions
    \item The Kullback-Leibler divergence measures how distinct the reward distribution $p_a$ is from the reward distribution given by choosing the best action $p_\star$.
    \begin{itemize}
        \item If $D_{KL}$ is high, that means that the actions we took are far from the optimal ones
        \item If, instead, $D_{KL}$ is low, that means that the distributions are very similar.
    \end{itemize}
    Note a thing: when the denominator $D_{KL}$ is small, the total regret goes up. This means you are forced to pull that suboptimal arm many times just to gather enough statistical evidence to prove it is actually worse than the best arm.
\end{itemize}

\subsection{Upper-Confidence-Bound action selection}
As we saw in section (\ref{lab:greedy_algorithms}), $\varepsilon$-greedy action selection forces the agent to explore, but it does so indiscriminately, with no preference for nearly-greedy actions or particularly uncertain actions. It would be better to favor non-greedy actions according to their potential for actually being optimal, taking into account how close their estimates are to being maximal and the uncertainties in those estimates. One effective way of doing this is to select actions according to:

\[
    A_t \doteq \argmax_{a} \left[ Q_t(a) + c \sqrt{\frac{\ln t}{N_t(a)}}\right], \qquad N_t(a)=\sum_i \mathbbm{1}(A_i=a)
\]

\noindent
The motifs behind the choice of this specific functional form are:
\begin{itemize}
    \item The formula explicitly splits the decision into \textit{exploitation} ($Q_t(a)$) and \textit{exploration} (the term under the square root)
    \item The term under the square root serves as an estimate of the variance. Each time an action $a$ is selected, in fact, the $N_t(a)$ term grows, lowering the overall uncertainty.
    \item Each time an action $a' \neq a$ is chosen, $t$ increases but $N_t(a)$ does not, and this makes the global uncertainty grow 
\end{itemize}

\begin{theorem}[Auer, Cesa-Bianchi, and Fischer, 2002]
    For the UCB algorithm, using $c=\sqrt{2}$, the expected total regret after $n$ steps is at most:
    \[
        L_t \leq \left[ 8 \sum_{a \mid \Delta_a > 0} \frac{\ln t}{\Delta_a} \right] + \left( 1 + \frac{\pi^2}{3}\right) \left( \sum_{j=1}^{K} \Delta_j \right)
     \]
\end{theorem}

Where $K$ is the number of machines the UCB policy is being run on.

This result is fundamental because it provides a \textbf{finite-time bound} on regret. While the Lai \& Robbins theorem was an asymptotic result (valid as $t \to \infty$), Auer et al. proved that UCB achieves \textbf{logarithmic regret} $\forall n > 0$. This confirms that UCB is not just a heuristic, but a mathematically optimal way to balance exploration and exploitation.